% Editable source for AnthonyWang_CV.pdf
% Compile with: pdflatex AnthonyWang_CV.tex
\documentclass[11pt,a4paper]{article}
\usepackage[margin=0.6in,top=0.55in,bottom=0.6in]{geometry}
\usepackage[T1]{fontenc}
\usepackage{newtxtext}
\usepackage{enumitem}
\usepackage[hidelinks]{hyperref}
\pagestyle{empty}
\setlength{\parindent}{0pt}
\setlength{\parskip}{0pt}
\setlist[itemize]{leftmargin=1.25em, label=\textbullet, nosep, topsep=0pt}

% Section and project formatting can be adjusted here.
\newcommand{\cvsection}[1]{\vspace{0.8em}\textbf{\MakeUppercase{#1}}\par\vspace{-0.5em}\noindent\rule{\linewidth}{0.4pt}\par\vspace{-0.5em}}
\newcommand{\project}[3]{\vspace{0.8em}\textbf{#1}\\\textit{#2}\\\textit{#3}\par}

\begin{document}
\begin{center}
    {\large\bfseries Yutong Wang}\\[2pt]
    +86 18912988997 \;|\; \href{mailto:anthonywang0705@gmail.com}{anthonywang0705@gmail.com} \;|\; \href{https://yutongwang777.github.io/}{https://yutongwang777.github.io/}
\end{center}
\vspace{-1em}

\cvsection{Education}
\textbf{South China University of Technology} \hfill Sep. 2022 - June 2026(expected)\\
B.Eng. Biomedical Engineering

\cvsection{Research Experience}
\textbf{Neural Field-Based Registration on Histopathology Images}\\
\textit{Xie Lab (University of California, Irvine)}\\
\textit{Research Intern (Supervisor: Prof. Xiaohui Xie)}
\begin{itemize}
    \item We proposed a neural field-based registration pipeline for histopathological images, enabling spatial alignment across different modalities of DAPI-stained tissues.
    \item Following a coarse-to-fine strategy, we first derive outline-based and patch-based affine transform matrix, and finally carry out a strongly supervised fusion step with the aid of Neural Field.
    \item We incorporated pretrained Cellpose-SAM into the pipeline. By jointly optimizing the registration loss and segmentation loss, the method improves registration speed while maintaining alignment accuracy.
\end{itemize}

\project{An Electron-Density Point-Cloud Framework for DPI Prediction\\(Under Review: \textit{Nature Communications})}{Xu Lab (SCUT)}{Core Member (Supervisor: Prof. Shidang Xu)}
\begin{itemize}
    \item Introduced, for the first time, electron cloud structures of molecules computed by quantum chemical computation methods into Drug--Protein Interaction prediction tasks.
    \item Developed E-CloudBind, a structure-based DPI model that fuses electron-density point clouds with covalent molecular graphs and a heterogeneous GNN to predict protein--ligand affinity.
    \item Created a density-adaptive interaction definition, yielding robustness to coordinate noise, low-resolution structures, and AlphaFold inputs.
    \item Achieved state-of-the-art accuracy on PDBbind datasets while maintaining minimal sensitivity to resolution/source; demonstrated superior generalization on out-of-distribution splits.
\end{itemize}

\project{Minimal High-Resolution Patches Are Sufficient for Whole Slide Image Representation via Cascaded Dual-Scale Reconstruction}{Xu Lab (SCUT)}{Contributing Member (Supervisor: Prof. Shidang Xu)}
\begin{itemize}
    \item We propose a two-stage sampling method based on attention mechanisms and clustering, whose sampling performance surpasses existing purely random sampling and single-stage sampling approaches.
    \item We design a self-supervised feature extractor for high-resolution images that, while drastically reducing the number of training samples, improves both training speed and model performance.
    \item On a classification task, the model uses only 4.5\% of the data yet improves over state-of-the-art methods by 6.3 percentage points in accuracy and 5.5 points in AUC.
\end{itemize}

\project{Concentric Dual Fusion Attention-MIL framework for WSI image segmentation}{Xu Lab (SCUT)}{Contributing Member (Supervisor: Prof. Shidang Xu)}
\begin{itemize}
    \item We propose a point-to-area feature enhancement method that leverages a channel-wise point-to-area attention mechanism to assign feature weights to each patch, enhancing representational capacity without requiring self-supervised pretraining.
    \item A point-to-point spatial fusion strategy is designed, employing a concentric-row attention mechanism to precisely control the influence of distant patches during feature fusion.
    \item The model demonstrates significantly improved performance over traditional MIL methods in multi-cancer subtype classification tasks, indicating stronger classification capability.
\end{itemize}

\end{document}
